\section{What must a computation retain?}
Eigen starts from a question about description: which distinctions must a computational system preserve to do useful work? A smaller representation is not necessarily a better one. It must still answer the question, predict the relevant change or reconstruct the information on which an action depends. We make that requirement explicit, account for its cost, and use the failure of an inadequate description to decide what to retain next.

Two tasks sharing one register provide a simple entry point. Sharing saves a resource, but the tasks may prefer different states. Their compromise has a price that further optimisation cannot remove. A summary of an evolving process faces a related choice: retain another variable, retain some history, or accept a specified prediction error. In a collection of records, the decisive distinction may be smaller still: whether a purchase preceded an approval. Keeping thousands of nearby passages is little help if the representation loses that relation.

The most useful evidence comes from constructions where we can see the choice change the answer. A conditional displacement absorbs a spin's coupling to its surrounding oscillators exactly, leaving a bounded correction from which its cold state can be predicted. A residual identifies a missing direction in a quantum trial space, and its excitation energy helps decide whether adding it is worthwhile. A compact continuation rule answers exact joint queries without scanning every record. These examples share a method: preserve the structure that makes the requested operation possible, then measure the cost of the remaining approximation.

A \emph{task} specifies the quantity to predict or operation to perform. A \emph{representation} is the object used in place of a more detailed description. Its \emph{retained distinctions} are the differences between situations that it can still express. A \emph{reconstruction} says what its smaller coordinates mean in the original space. A \emph{residual} measures the failure of a specified calculation to commute with that reduction. These definitions give compression a purpose. Total cost includes constructing, storing, updating and using the representation, together with recovering witnesses or correcting an inadequate answer.

Eigen's contribution is a proposed synthesis and a method of construction, developed across mathematical papers, finite computational prototypes and a working research environment. Familiar objects - probability distributions, projections, effective operators and dependency graphs - acquire particular jobs within that method. The claim is strongest where the maps, assumptions and comparison can be written down. The larger proposition is that treating representation choice as part of the calculation can expose useful structure that a single global description obscures.

The article follows that proposition through three questions. First, when does an additional distinction earn its cost? The shared-register example leads to predictive loss, coarse dynamics and effective physical laws. Second, how does the account change a practical procedure? Quantum trial-space repair, spectral filtering and Dynamic Memory provide distinct answers. Third, how can the resulting knowledge remain usable when its sources or assumptions change? The Writing App retains scientific objects, their grounds and authorised transitions, providing an implemented counterpart and a route towards an extensible software development kit (SDK).

The evidence has different strengths at these three levels. Finite mathematical identities explain mechanisms. Simulations and archived comparisons establish particular outcomes under stated conditions. The software demonstrates persistent objects and recoverable operations. The SDK describes a product direction to evaluate. We develop each at the point where it advances the argument, so the reader can judge both the present contribution and the next useful test.
